\documentclass[10pt,a4paper]{article}
\usepackage[margin=1.8cm]{geometry}
\usepackage[utf8]{inputenc}\usepackage[LGR,T1]{fontenc}\usepackage{lmodern}\usepackage{textalpha}\usepackage{substitutefont}
\substitutefont{LGR}{\rmdefault}{cmr}\substitutefont{LGR}{\sfdefault}{cmss}\substitutefont{LGR}{\ttdefault}{cmtt}
\usepackage{booktabs,array,enumitem,xcolor}\usepackage[hidelinks]{hyperref}
\sloppy\emergencystretch=3em\setlength{\parindent}{0pt}\setlength{\parskip}{5pt}
\setlist{nosep,leftmargin=*}\renewcommand{\arraystretch}{1.22}
\DeclareUnicodeCharacter{2212}{\ensuremath{-}}
\newcommand{\src}[1]{\par{\scriptsize\color{gray}Πηγές: {\ttfamily #1}}\par}
\newcommand{\page}[1]{\clearpage\section*{#1}}
\begin{document}
\begin{center}{\LARGE Οι στόχοι σου και τα πειράματα}\\[5pt]{\large QRL-Lab · Αναφορά προς τον Κώστα · 10 Νοεμβρίου 2026}\end{center}
\textbf{Πού βρισκόμαστε.} Έχουμε συγκρίσεις κλασικών μοντέλων σε προσομοίωση και ολοκληρωμένο τελικό έλεγχο στο βασικό φορτίο αιτημάτων. Δεν έχει επιβεβαιωθεί ότι το Dinos+Transfer V.2 είναι συνολικά το καλύτερο. Η τελευταία παρτίδα ολοκληρώθηκε: τα νέα αποτελέσματά της περιλαμβάνονται στη σελίδα 7. Οι κβαντικές δοκιμές και η αξιολόγηση στο πραγματικό σύστημα παραμένουν σε παύση.
\src{docs/\allowbreak STATE.md §§1-\allowbreak 5; M-\allowbreak 20261109-\allowbreak 040; M-\allowbreak 20261110-\allowbreak 017; .lab/\allowbreak j450/\allowbreak T-\allowbreak summary.json}
\textbf{Άνοιγμα αρχείου:} στην καρτέλα \textbf{Αρχεία} (εικονίδιο φακέλου), γράψε τη διαδρομή\\
{\small\texttt{\detokenize{~/qrl-lab/reports/2026-11-10-goals-experiments/report.pdf}}}
\subsection*{Τι έχεις ζητήσει και σε ποια κατάσταση βρίσκεται}
\begin{tabular}{p{10.6cm}p{4.6cm}}\toprule
Στόχος & Κατάσταση\\\midrule
Δημοσίευση για υβριδικό κβαντικό--κλασικό σύστημα που αυξομειώνει τους πόρους εφαρμογών. Πρώτα προσομοίωση, μετά πραγματικό Kubernetes/Redis. & Σε εξέλιξη το κλασικό σκέλος· κβαντικό και πραγματικό σύστημα σε παύση.\\
Περίληψη για CCGrid έως 24/11/2026· αρχικός στόχος υποβολής άρθρου μέσα/τέλος Νοεμβρίου. & Σε εξέλιξη· υποβολή δεν τεκμηριώνεται.\\
Καλύτερη πρόβλεψη φορτίου, ταχύτερη μάθηση, μικρότερο κόστος χρόνου/μνήμης/δίσκου, καλύτερη προσαρμογή. Επιλογή 2--3 κατάλληλων κβαντικών προσεγγίσεων. & Μερική κλασική διερεύνηση· ολοκληρωμένη απόδειξη αυτών των πλεονεκτημάτων δεν υπάρχει. Κβαντικές δοκιμές σε παύση.\\
Πιο ρεαλιστικά αιτήματα τύπου Poisson, απότομες αιχμές, σήματα καθυστέρησης και χρήσης πόρων. & Ολοκληρώθηκε η υλοποίηση και έγιναν συγκρίσεις· έλεγχος σε πραγματικό σύστημα εκκρεμεί.\\
Σύγκριση όλων των μοντέλων και προσαρμοσμένου DRe-SCale, με κοινές συνθήκες και επαναλήψεις. & Ολοκληρώθηκε στο βασικό φορτίο· οι νεότερες αιχμές παραμένουν σε διερευνητική αξιολόγηση.\\
Σαφείς ενημερώσεις και προσβάσιμα αποτελέσματα από κινητό, χωρίς εσωτερική ορολογία. & Σε εξέλιξη· το αίτημά σου ανανεώθηκε σήμερα.\\\bottomrule
\end{tabular}
\src{docs/\allowbreak PLAN.md §1, R4/\allowbreak R6/\allowbreak R7/\allowbreak R8/\allowbreak R10; docs/\allowbreak researcher\_requirements.md A1-\allowbreak A7, B8, D4, F1-\allowbreak F2; docs/\allowbreak STATE.md §§1-\allowbreak 3; docs/\allowbreak rules/\allowbreak paused.md; M-\allowbreak 20261110-\allowbreak 018}
\subsection*{Πώς διαβάζονται οι αριθμοί}
Η \textbf{βαθμολογία} συνδυάζει εξυπηρετούμενα αιτήματα, κόστος πόρων, καθυστερήσεις και απορρίψεις· μεγαλύτερη είναι καλύτερη στο ίδιο περιβάλλον. Δεν είναι ποσοστό ακρίβειας. \textbf{Έλεγχος ανάπτυξης} σημαίνει δεδομένα που χρησιμοποιούμε όσο επιλέγουμε ρυθμίσεις. \textbf{Τελικός έλεγχος} σημαίνει ξεχωριστά δεδομένα που εξετάζονται μία φορά. Αριθμοί διαφορετικών περιβαλλόντων δεν συγκρίνονται απευθείας.
\src{docs/\allowbreak cycle2\_plan.md §§2,4,6; docs/\allowbreak paper\_notes.md §3}
\page{2. Οι παλιότερες δοκιμές: τι μάθαμε}
Οι παρακάτω οικογένειες καλύπτουν τις αρχικές παρτίδες και την ανάπτυξη πριν από τους τελικούς ελέγχους. Είναι ιστορικό της έρευνας· τα διερευνητικά ευρήματα δεν γίνονται αυτόματα τελικά συμπεράσματα.
\begin{tabular}{p{3.0cm}p{6.2cm}p{5.8cm}}\toprule
Δοκιμή / σκοπός & Συνθήκες και αποτέλεσμα & Κατάσταση και όριο\\\midrule
Αρχική επαλήθευση και αναπαραγωγή, 16/10 & Η σταθερή ελάχιστη υποδομή ξεπερνούσε αποθηκευμένα μοντέλα στον παλιό προσομοιωτή. Αναπαραγωγή αποθηκευμένων βαθμολογιών περίπου εντός 1\%. & Ολοκληρώθηκε. Η καθυστέρηση και η χρήση πόρων δεν εξαρτιόνταν σωστά από τις ενέργειες σε εκείνα τα αρχεία.\\
Πρώτες προσομοιώσεις W0, 17--19/10 & Έλεγχοι αναφοράς, ρύθμιση και μεταφορά μάθησης. Η αρχική εξερευνητική βαθμολογία του DERP ήταν αρνητική. & Ολοκληρώθηκαν/σταμάτησαν επιμέρους δοκιμές. Αξιόπιστη συνολική αποτίμηση ρυθμίσεων: άγνωστη στο μητρώο. Το παλιό κριτήριο βαθμολόγησης δεν στηρίζει νέο ισχυρισμό.\\
Πραγματικό Kubernetes, στάδια Α/Β, 16--18/10 & Προετοιμασία εφαρμογής και αύξηση ταυτόχρονων αιτημάτων. Προβλήματα μνήμης και μετακίνησης δεδομένων σταμάτησαν τη βαθμονόμηση. & Σε παύση. Δεν ολοκληρώθηκε σύγκριση των τελικών μοντέλων στο πραγματικό σύστημα. Η χωριστή αλλαγή ρυθμίσεων 16/10 καταγράφεται χωρίς τεκμηριωμένο ερευνητικό αποτέλεσμα.\\
Πρώτες αλλαγές προσομοιωτή, κύκλοι 1--8 & Νέο περιβάλλον, αντιγραφή παλιού, συνδεδεμένα σήματα πόρων, διόρθωση μνήμης LSTM και κλίμακας εισόδων. & Ολοκληρώθηκαν. Κρατήθηκαν διορθώσεις μάθησης. Η μεγάλη παλιά αναλογία περίπου 17 φορές εντοπίστηκε ως πρόβλημα πρωτοκόλλου σύγκρισης.\\
Αρχικό φορτίο και καθυστέρηση εκκίνησης, κύκλοι 9--13 & Δοκιμές ημιτονοειδούς φορτίου/αιχμών και χρόνου εκκίνησης πόρων. Μεγάλες καθυστερήσεις προκάλεσαν αστάθεια. & Ολοκληρώθηκαν. Επιλέχθηκε βασικό φορτίο με αιχμή και μικρή καθυστέρηση· η επιλογή έγινε κατά την ανάπτυξη.\\
Αντοχή σε αλλαγές, κύκλοι 14--17 & Περισσότερες αιχμές, θόρυβος και καθυστέρηση. Δεν διατηρήθηκε παντού η επιθυμητή κατάταξη. & Ολοκληρώθηκαν. Το επιτυχημένο εύρημα περιοριζόταν σε συγκεκριμένες συνθήκες.\\
Αρχιτεκτονικές και μέθοδοι μάθησης, κύκλοι 18--20 & Σύγκριση αναδρομικών/πυκνών δικτύων, προεκπαίδευσης και ενημερώσεων με περισσότερα βήματα. Κρατήθηκε η αρχιτεκτονική Dinos με ιστορικό. & Ολοκληρώθηκαν. Κάποιες βελτιώσεις ανάπτυξης δεν επαναλήφθηκαν στον τελικό έλεγχο.\\
Μεταγενέστερη διερεύνηση, κύκλοι 22--23 & Έλεγχος διακύμανσης μεταξύ επαναλήψεων και πακέτου αλλαγών. Το πακέτο δεν υιοθετήθηκε. & Ολοκληρώθηκαν. Χωριστή πλήρης σύνοψη της μελέτης διακύμανσης: άγνωστη στο μητρώο.\\\bottomrule
\end{tabular}
\src{docs/\allowbreak EXPERIMENTS.md: αρχικές παρτίδες R-\allowbreak 20261016-\allowbreak 011/\allowbreak -\allowbreak 012/\allowbreak -\allowbreak 020/\allowbreak -\allowbreak 021/\allowbreak -\allowbreak 022, R-\allowbreak 20261017-\allowbreak 004/\allowbreak -\allowbreak 010/\allowbreak -\allowbreak 013/\allowbreak -\allowbreak 015/\allowbreak -\allowbreak 021/\allowbreak -\allowbreak 022 και κύκλοι c1-\allowbreak c24; docs/\allowbreak tuning\_experiments.md κύκλοι1-\allowbreak 21; docs/\allowbreak PLAN.md §2 F1; .lab/\allowbreak tuning/\allowbreak cycles/\allowbreak c23\_screen.md}
Το παλιό σύστημα HPA αντιδρούσε σε τυχαία σήματα πόρων, άρα η μεγάλη απόσταση από αυτό δεν αρκεί. Στις νεότερες συγκρίσεις συμπεριλαμβάνεται και \textbf{ρυθμισμένος ελεγκτής κανόνων}, που προσαρμόζει τους πόρους στη ζήτηση. Αυτός είναι ουσιαστικά δυσκολότερος αντίπαλος.
\src{docs/\allowbreak tuning\_experiments.md Notes; docs/\allowbreak paper\_notes.md §§3,5}
\page{3. Οι αριθμητικοί στόχοι και οι τελικοί έλεγχοι του πρώτου κύκλου}
\textbf{Τι εννοούσες με «τα νούμερα που θέλουμε».} Το μήνυμά σου στις 2/11 επαναφέρει τον στόχο, χωρίς να δίνει νέο αριθμητικό όριο. Ο παλιός δείκτης προόδου είχε ως σημεία αναφοράς συνολική βαθμολογία 3.800.000 για DInos και 220.000 για DERP (λόγος περίπου 17,3). Είναι αθροίσματα του παλιού πρωτοκόλλου, όχι στόχοι ανά επεισόδιο για τα νεότερα πειράματα. Δεν βρέθηκε εγκεκριμένο νέο απόλυτο όριο για το φορτίο Poisson: \textbf{άγνωστο}.
\src{M-\allowbreak 20261102-\allowbreak 004; code-\allowbreak poisson-\allowbreak spiky/\allowbreak tuning/\allowbreak closeness.py:22-\allowbreak 31; docs/\allowbreak paper\_notes.md §5; M-\allowbreak 20261029-\allowbreak 017}
Στις 22/10 αναγνώρισες ότι οι παλιοί απόλυτοι αριθμοί δεν μπορούν να ζητούνται αυτούσιοι μετά τις διορθώσεις. Για τον δεύτερο κύκλο έχει καταγραφεί ως επιθυμητή σειρά:
\begin{center}\small Dinos+Transfer V.2 $>$ LSTM+transfer $>$ Dinos $>$ LSTM\\ $>$ DRe-SCale $>$ DERP $>$ HPA.\end{center}
Η σειρά είναι στόχος που ελέγχεται, όχι αποτέλεσμα που θεωρείται δεδομένο. Η αλήθεια των μετρήσεων προηγείται της επιθυμητής κατάταξης.
\src{M-\allowbreak 20261022-\allowbreak 050; docs/\allowbreak researcher\_requirements.md A3-\allowbreak A5, H1 και απάντηση R-\allowbreak 20261030-\allowbreak 006; docs/\allowbreak cycle2\_plan.md §1}
\subsection*{Δύο ανεξάρτητοι τελικοί έλεγχοι στο παλιότερο φορτίο αιχμής}
Κάθε έλεγχος είχε πέντε νέες επαναλήψεις εκπαίδευσης και σαράντα έξι καταγεγραμμένες εκτελέσεις μαζί με τις αναφορές. Τα δεδομένα τους έχουν ήδη χρησιμοποιηθεί και δεν ξαναδιαβάστηκαν για αυτή την αναφορά. Οι αριθμοί προέρχονται από τις υπάρχουσες περιλήψεις.
\begin{center}\begin{tabular}{lrr}\toprule
Μοντέλο & Πρώτος τελικός έλεγχος & Δεύτερος τελικός έλεγχος\\\midrule
LSTM+transfer &90.635&89.331\\
LSTM &89.966&87.347\\
Dinos &86.148&88.026\\
Dinos+Transfer V.2 &87.781&83.983\\
DERP με ιστορικό &69.803&68.592\\
DERP αρχικό &62.998&62.228\\
Ελεγκτής κανόνων βάσει ζήτησης &47.530&47.572\\\bottomrule
\end{tabular}\end{center}
\src{.lab/\allowbreak tuning/\allowbreak cycles/\allowbreak c21\_test10\_result.md; .lab/\allowbreak tuning/\allowbreak cycles/\allowbreak c24\_test11\_result.md; docs/\allowbreak EXPERIMENTS.md R-\allowbreak 20261022-\allowbreak 005 και R-\allowbreak 20261023-\allowbreak 001}
\textbf{Πρώτη υπόθεση του πρώτου κύκλου:} το Dinos+Transfer V.2 υπερτερεί του αρχικού DERP. Επιβεβαιώθηκε και στους δύο ελέγχους: διαφορές περίπου +24.784 και +21.755 μονάδες, αντίστοιχα.

\textbf{Δεύτερη υπόθεση του πρώτου κύκλου:} υπερτερεί του DERP με ιστορικό. Επιβεβαιώθηκε επίσης: +17.979 και +15.391 μονάδες.

Η επόμενη υπόθεση, ότι το Dinos υπερτερεί του LSTM, δεν επιβεβαιώθηκε. Σύμφωνα με τον προκαθορισμένο κανόνα, οι επόμενες συγκρίσεις έμειναν περιγραφικές. Δεν αποδείχθηκε ότι το Dinos+Transfer V.2 είναι πρώτο. Η νωρίτερη καλή επίδοση κατά την προσαρμογή δεν αποδεικνύει συνολικά ταχύτερη μάθηση, επειδή εκείνες οι καμπύλες δεν χρέωναν την προεκπαίδευση.
\src{.lab/\allowbreak tuning/\allowbreak cycles/\allowbreak c21\_test10\_result.md: Fixed-\allowbreak sequence και Reading; .lab/\allowbreak tuning/\allowbreak cycles/\allowbreak c24\_test11\_result.md: Fixed sequence και Reading}
\page{4. Πρώτες παρτίδες Poisson και αιχμών}
\textbf{Σκοπός.} Να ελεγχθούν οι μέθοδοι σε αφίξεις αιτημάτων που μεταβάλλονται τυχαία, με ουρά αναμονής, κόστος πόρων και ποινές για κακή εξυπηρέτηση. Οι παρακάτω παρτίδες χρησιμοποιούσαν δεδομένα ανάπτυξης· κανένα εύρημα εδώ δεν είναι νέα τελική επιβεβαίωση.
\begin{tabular}{p{3.5cm}p{7.4cm}p{4.1cm}}\toprule
Παρτίδα & Συνθήκες / κύριο αποτέλεσμα & Όριο\\\midrule
Αρχικοί έλεγχοι λειτουργίας και κόστους, 29/10 & Δύο ομάδες από 14 εκτελέσεις ολοκληρώθηκαν. Οι μετρήσεις κόστους επέβαλαν αλλαγή σχεδίου. & Δεν είναι σύγκριση τελικών μοντέλων.\\
Πρώτη αναζήτηση ρυθμίσεων, 29--30/10 & 48 εκτελέσεις ολοκληρώθηκαν μετά αποκατάσταση καταγραφής. Με 2 επαναλήψεις ανά επιλογή, Dinos περίπου 196,5 χιλ. και επιλεγμένο Dinos+Transfer V.2 περίπου 188,9 χιλ. & Πολύ λίγες επαναλήψεις· μεγάλη διαφορά μεταξύ ομάδων.\\
Επέκταση αναζήτησης, 31/10--1/11 & Προηγήθηκαν διακοπές από ελέγχους του προγράμματος. Στην ολοκληρωμένη επέκταση καταγράφηκαν 192 υποβολές. Καλύτερη μέση επίδοση Dinos+Transfer V.2 περίπου 200,7 χιλ. στην τελική φάση ανάπτυξης. & Διαφορετικές πεντάδες επαναλήψεων ανά ομάδα· η απόσταση δεν αποδίδεται μόνο στη μέθοδο.\\
Πρώτη εκπαίδευση με έντονες αιχμές, 1--2/11 & Ολοκληρώθηκαν 7 μοντέλα επί 6 επαναλήψεις. Μέσοι: DERP με ιστορικό 190,1 χιλ., Dinos+Transfer V.2 187,9, DERP 186,0, Dinos 184,6, LSTM+transfer 177,6, LSTM 165,3, DRe-SCale 57. & Διακοπή άλλων εκτελέσεων λόγω ορίου χρόνου. Μεγάλη αστάθεια DRe-SCale.\\
Συμπλήρωση και έλεγχος χρήσης ιστορικού, 2/11 & 24/24 εκτελέσεις: ανακατεμένο ιστορικό Dinos 178,5 χιλ., μόνο πρόσφατη είσοδος 162,0, πυκνό δίκτυο 159,7, πυκνό με προεκπαίδευση 154,3. & 6 επαναλήψεις· μεγάλες διακυμάνσεις. Δεν απομονώθηκε σταθερό όφελος από το ιστορικό.\\\bottomrule
\end{tabular}
\src{docs/\allowbreak EXPERIMENTS.md R-\allowbreak 20261029-\allowbreak 006/\allowbreak -\allowbreak 007/\allowbreak -\allowbreak 010, R-\allowbreak 20261031-\allowbreak 001/\allowbreak -\allowbreak 002, R-\allowbreak 20261101-\allowbreak 002/\allowbreak -\allowbreak 005, R-\allowbreak 20261102-\allowbreak 001; .lab/\allowbreak j253/\allowbreak summary.json; .lab/\allowbreak j289/\allowbreak summary.json; .lab/\allowbreak j299/\allowbreak stageA\_summary.json; .lab/\allowbreak j312/\allowbreak summary.json}
\subsection*{Προτάσεις που δεν είναι επιπλέον αποτελέσματα}
Οι ανεξάρτητες προτάσεις αναζήτησης του δεύτερου κύκλου, το προσχέδιο διευρυμένης επανάληψης στις 2/11 και το αρχικό σχέδιο νέας αναζήτησης στις 6/11 καταγράφονται στο μητρώο ως σχέδια ή προετοιμασίες. Δεν προστίθενται ως νέες μετρήσεις. Η αναζήτηση στις 6--9/11 προχώρησε αργότερα στις συγκεκριμένες παρτίδες της επόμενης σελίδας.
\src{docs/\allowbreak EXPERIMENTS.md: cycle2-\allowbreak codex J25, C330/\allowbreak J332, C341/\allowbreak J350, C345 W/\allowbreak F}
\subsection*{Τι σημαίνει «μεταφορά μάθησης» εδώ}
Συνήθως σημαίνει ότι ένα μοντέλο μαθαίνει πρώτα σε μια κατάσταση και μετά προσαρμόζεται σε άλλη. Οι νεότερες συνταγές του Dinos+Transfer V.2 στις αιχμές και στη μεταβολή φορτίου κάνουν δύο φάσεις μάθησης \emph{στο ίδιο είδος φορτίου}, με διαφορετική ροή τυχαίων γεγονότων. Επομένως, τα αποτελέσματά τους δεν αποδεικνύουν από μόνα τους μεταφορά σε διαφορετικό περιβάλλον. Επιπλέον, Dinos και Dinos+Transfer V.2 έχουν διαφορετικές ρυθμίσεις μάθησης· συγκρίνονται ολόκληρες συνταγές.
\src{.lab/\allowbreak j445/\allowbreak FREEZE.md: Ερώτημα και περιορισμός, Πίνακας εκτελέσεων; reports/\allowbreak 2026-\allowbreak 11-\allowbreak 10-\allowbreak models-\allowbreak report/\allowbreak report.tex: LSTM+transfer, Dinos+Transfer V.2}
\page{5. Βασικό φορτίο Poisson: αναζήτηση, επανάληψη, τελικός έλεγχος}
\textbf{Αναζήτηση ρυθμίσεων, 7/11.} Ολοκληρώθηκαν 55/55 εργασίες με δύο επαναλήψεις ανά υποψήφια ρύθμιση. Σκοπός ήταν η επιλογή συνταγών, όχι ο τελικός ισχυρισμός υπεροχής.

\textbf{Επανάληψη επιλεγμένων συνταγών, 7--8/11.} Δέκα μοντέλα, δέκα επαναλήψεις εκπαίδευσης ανά μοντέλο, συνολικά 100 τελικά αρχεία αποτελεσμάτων. Κοινό συνολικό όριο 39.936 βημάτων εκπαίδευσης και 20 διαδρομές ανάπτυξης. Παλιές εκτελέσεις που αντικαταστάθηκαν δεν μετρήθηκαν δεύτερη φορά.

\textbf{Τελικός έλεγχος, 9/11.} Ολοκληρώθηκαν 103 μονάδες αξιολόγησης: τα αποθηκευμένα μοντέλα και οι τρεις ελεγκτές αναφοράς, σε 20 νέες κοινές διαδρομές. Προηγήθηκε επιβεβαίωση ότι η φόρτωση των μοντέλων αναπαρήγαγε τα παλιά αποτελέσματα. Αρχικές τεχνικές αποτυχίες διορθώθηκαν πριν ολοκληρωθεί η προβλεπόμενη αξιολόγηση. Δεν επανεξετάστηκαν τα δεδομένα τελικού ελέγχου για αυτή την αναφορά.
\src{docs/\allowbreak EXPERIMENTS.md: C345 W screening J380, F J408, integrity J410, test12 J414/\allowbreak J419/\allowbreak J421; docs/\allowbreak STATE.md §3; M-\allowbreak 20261109-\allowbreak 040; .lab/\allowbreak c401/\allowbreak RESULTS.md}
\begin{center}\begin{tabular}{lrrr}\toprule
Μοντέλο & Ανάπτυξη & Τελικός έλεγχος & Παραβίαση εξυπηρέτησης\\\midrule
Dinos &200&202&7\%\\
Πυκνό δίκτυο ελέγχου &197&199&5\%\\
DERP με διορθωμένη μάθηση &195&197&7\%\\
Πυκνό δίκτυο με προεκπαίδευση &192&194&9\%\\
Dinos+Transfer V.2 &189&190&19\%\\
LSTM+transfer &186&188&9\%\\
LSTM &186&188&11\%\\
DERP με ιστορικό &184&186&12\%\\
DERP αρχικό &131&132&9\%\\
DRe-SCale προσαρμοσμένο &−17&−19&63\%\\\midrule
Αρχικός HPA &131&133&2\%\\
Ρυθμισμένος ελεγκτής κανόνων &212&214&5\%\\
Σταθεροί πόροι &135&138&11\%\\\bottomrule
\end{tabular}\end{center}
Οι βαθμολογίες είναι σε χιλιάδες μονάδες ανά επεισόδιο 5.000 βημάτων, στρογγυλοποιημένες όπως στις υπάρχουσες δημοσιευμένες περιλήψεις. Η τελευταία στήλη αφορά τον τελικό έλεγχο: ποσοστό βημάτων όπου δεν τηρείται ο στόχος ποιότητας εξυπηρέτησης στον προσομοιωτή.
\src{.lab/\allowbreak c401/\allowbreak RESULTS.md; M-\allowbreak 20261109-\allowbreak 040}
\textbf{Πρώτη υπόθεση του δεύτερου κύκλου:} το Dinos+Transfer V.2 υπερτερεί του προσαρμοσμένου DRe-SCale. Υποστηρίχθηκε: διαφορά +209,5 χιλ., με διάστημα αβεβαιότητας 95\% από +140,8 έως +277,5 χιλ.

\textbf{Δεύτερη υπόθεση:} υπερτερεί του DERP με ιστορικό. Δεν υποστηρίχθηκε: +4,8 χιλ., με διάστημα από −8,6 έως +20,5 χιλ. Οι επόμενες υποθέσεις δεν εξετάστηκαν ως επιβεβαιωτικές. Ο ρυθμισμένος ελεγκτής κανόνων ξεπερνά όλα τα εκπαιδευμένα μοντέλα σε μέση βαθμολογία. Η κακή επίδοση της δικής μας προσαρμογής DRe-SCale δεν αποδεικνύει ότι η πρωτότυπη δημοσιευμένη μέθοδος είναι γενικά κακή ή ότι αναπαράχθηκε πλήρως.
\src{M-\allowbreak 20261109-\allowbreak 040; docs/\allowbreak cycle2\_plan.md §§3,6; docs/\allowbreak paper\_notes.md §4}
\page{6. Νεότερες συνθήκες: αιχμές και αλλαγή φορτίου}
\subsection*{Εκπαίδευση με απότομες αιχμές, 9--10/11}
Επτά μοντέλα επί δέκα επαναλήψεις: 70/70 επιτυχείς εκτελέσεις. Κάθε συνταγή είχε 19.968 συνολικά βήματα εκπαίδευσης, μαζί με την προεκπαίδευση όπου υπήρχε, και αξιολογήθηκε σε 20 κοινές διαδρομές ανάπτυξης. Σκοπός ήταν να φανεί αν οι προηγούμενες ρυθμίσεις αντέχουν σε πιο απότομα φορτία, χωρίς νέα αναζήτηση ρυθμίσεων.
\begin{center}\begin{tabular}{lrrr}\toprule
Μοντέλο & Μέση βαθμολογία & Τυπική απόκλιση & Παραβίαση εξυπηρέτησης\\\midrule
Dinos & 191,6 & $\pm$7,1 & 10,6\% \\
Dinos+Transfer V.2 & 188,1 & $\pm$12,0 & 10,3\% \\
LSTM+transfer & 179,9 & $\pm$7,7 & 14,5\% \\
LSTM & 164,6 & $\pm$10,9 & 22,4\% \\
DERP αρχικό & 153,6 & $\pm$29,4 & 4,8\% \\
DERP με ιστορικό & 91,5 & $\pm$149,8 & 35,5\% \\
DRe-SCale προσαρμοσμένο & 86,0 & $\pm$113,9 & 38,6\% \\
\bottomrule\end{tabular}\end{center}
Βαθμολογίες και τυπική απόκλιση σε χιλιάδες μονάδες. Η τυπική απόκλιση δείχνει πόσο αλλάζει το αποτέλεσμα όταν επαναλαμβάνεται η εκπαίδευση· δεν είναι διάστημα βεβαιότητας. Αναφορές στο ίδιο φορτίο: αρχικός HPA 152,4 χιλ., ρυθμισμένος ελεγκτής 201,9 χιλ., ελεγκτής με γνώση μελλοντικού φορτίου 222,7 χιλ. Οι ρυθμίσεις των αναφορών επιλέχθηκαν στα ίδια δεδομένα ανάπτυξης.
\src{.lab/\allowbreak j442/\allowbreak summary.json; docs/\allowbreak EXPERIMENTS.md batch S; code-\allowbreak poisson-\allowbreak spiky/\allowbreak tuning/\allowbreak results/\allowbreak spiky-\allowbreak stage1/\allowbreak spiky1-\allowbreak P0-\allowbreak reference-\allowbreak off-\allowbreak s16000.json}
Το Dinos είναι το καλύτερο εκπαιδευμένο μοντέλο κατά μέσο όρο, όμως κάτω από τον ρυθμισμένο ελεγκτή κανόνων. Η διαφορά του Dinos+Transfer V.2 από το Dinos είναι −3,5 χιλ. με μεγάλη διασπορά 16,9 χιλ.· δεν τεκμηριώνεται καθαρή υπεροχή του. DERP με ιστορικό και DRe-SCale παρουσιάζουν σοβαρή αστάθεια.
\src{.lab/\allowbreak j442/\allowbreak summary.json; docs/\allowbreak EXPERIMENTS.md J442}
\subsection*{Αλλαγή συνθηκών χωρίς νέα εκπαίδευση, 9/11}
Αξιολογήθηκαν 90 αποθηκευμένα μοντέλα από το βασικό φορτίο: εννέα μέθοδοι επί δέκα επαναλήψεις, σε δέκα διαδρομές με διαφορετικές αιχμές και περισσότερο θόρυβο. Dinos: 211,7 χιλ. με τυπική απόκλιση 12,7 χιλ. Αναφορές: αρχικός HPA 166,7 χιλ., ρυθμισμένος ελεγκτής 225,0 χιλ., ελεγκτής με γνώση του μέλλοντος 238,6 χιλ. Κανένα μοντέλο δεν ξεπέρασε τον ρυθμισμένο ελεγκτή. Πρόκειται για δεδομένα ανάπτυξης, χωρίς προσαρμογή των βαρών σε αυτή την παρτίδα.
\src{.lab/\allowbreak j438/\allowbreak summary.json; .lab/\allowbreak j431/\allowbreak REPORT.md; docs/\allowbreak EXPERIMENTS.md c3-\allowbreak shift-\allowbreak reference/\allowbreak c3-\allowbreak shift-\allowbreak eval}
\subsection*{Γραφικές συμπεριφοράς, 10/11}
Ολοκληρώθηκαν οκτώ πρόσθετες αξιολογήσεις για την προηγούμενη αναφορά: έξι αποθηκευμένα μοντέλα και δύο ομάδες ελεγκτών αναφοράς. Οι καμπύλες πόρων/φορτίου χρησιμοποιούν μία επανάληψη εκπαίδευσης ανά μοντέλο και τρεις διαδρομές ανάπτυξης· δεν είναι μέσος όρος όλων των επαναλήψεων. Οι βαθμολογίες αναπαράχθηκαν ακριβώς. Λείπει η καμπύλη συμπεριφοράς του DRe-SCale. Οι σκόπιμες αποτυχίες ελέγχων ασφαλούς λειτουργίας καταγράφονται χωριστά, δεν είναι αποτυχημένα εκπαιδευμένα μοντέλα.
\src{.lab/\allowbreak j446/\allowbreak REPORT.md; docs/\allowbreak EXPERIMENTS.md J444/\allowbreak J446; reports/\allowbreak 2026-\allowbreak 11-\allowbreak 10-\allowbreak models-\allowbreak report/\allowbreak report.pdf}
Ο ελεγκτής που γνωρίζει το μελλοντικό φορτίο είναι βοηθητικό σημείο σύγκρισης, \textbf{όχι αποδεδειγμένο βέλτιστο όριο} και όχι εφαρμόσιμος ανταγωνιστής χωρίς αυτή τη γνώση.
\src{docs/\allowbreak paper\_notes.md §6}
\page{7. Η τελευταία ολοκληρωμένη παρτίδα: νέα επανάληψη}
Σκοπός: να συγκριθούν ξανά το Dinos και το Dinos+Transfer V.2 με νέες επαναλήψεις εκπαίδευσης. Δύο είδη φορτίου, δέκα επαναλήψεις ανά μοντέλο και φορτίο: \textbf{40/40 επιτυχείς εκτελέσεις, χωρίς επανάληψη αποτυχημένης προσπάθειας}. Η συγκεκριμένη παρτίδα έχει τελειώσει· δεν εκτελέστηκε κανένα νέο πείραμα για το παρόν PDF.

Κοινό συνολικό κόστος αλληλεπίδρασης: 19.968 βήματα εκπαίδευσης. Dinos: μία φάση. Dinos+Transfer V.2: 9.984 βήματα προεκπαίδευσης και 9.984 προσαρμογής στο ίδιο είδος φορτίου. Ένα προκαθορισμένο τελικό μοντέλο ανά εκτέλεση· αξιολόγηση σε 20 κοινές διαδρομές με αιχμές και 10 με αλλαγή συνθηκών.
\src{.lab/\allowbreak j445/\allowbreak FREEZE.md; .lab/\allowbreak j450/\allowbreak summarize\_t.py; .lab/\allowbreak j450/\allowbreak T-\allowbreak summary.json; .lab/\allowbreak j450/\allowbreak launcher-\allowbreak summary.txt; M-\allowbreak 20261110-\allowbreak 017}
\begin{center}\begin{tabular}{llrrr}\toprule
Φορτίο & Μοντέλο & Βαθμολογία & Τυπική απόκλιση & Παραβιάσεις\\\midrule
Αιχμές & Dinos & 185,8 & $\pm$18,6 & 15,2\% \\
Αιχμές & Dinos+Transfer V.2 & 185,5 & $\pm$10,4 & 14,5\% \\
Αλλαγή συνθηκών & Dinos & 215,5 & $\pm$7,4 & 17,2\% \\
Αλλαγή συνθηκών & Dinos+Transfer V.2 & 200,5 & $\pm$17,5 & 13,5\% \\
\bottomrule\end{tabular}\end{center}
Βαθμολογία και απόκλιση σε χιλιάδες μονάδες ανά επεισόδιο. Παραβιάσεις: ποσοστό βημάτων χωρίς την απαιτούμενη ποιότητα εξυπηρέτησης. Και εδώ χρησιμοποιούνται μόνο δεδομένα ανάπτυξης.
\subsection*{Τι δείχνει η διαφορά των δύο μεθόδων}
\textbf{Στις αιχμές}: Dinos+Transfer V.2 μείον Dinos = -0,3 χιλ. Το διάστημα αβεβαιότητας 95\% είναι [-10,2 έως 12,1] χιλ.

\textbf{Στην αλλαγή συνθηκών}: Dinos+Transfer V.2 μείον Dinos = -15,0 χιλ. Το διάστημα αβεβαιότητας 95\% είναι [-27,6 έως -4,1] χιλ.

Στις αιχμές η διαφορά είναι πρακτικά μηδενική σε σχέση με την αβεβαιότητα. Στην αλλαγή συνθηκών το Dinos έχει υψηλότερη βαθμολογία σε αυτή την επανάληψη. Το Dinos+Transfer V.2 έχει λιγότερες παραβιάσεις κατά μέσο όρο, χωρίς χωριστό στατιστικό έλεγχο αυτής της διαφοράς, και χρησιμοποιεί περισσότερες ενεργές μονάδες εφαρμογής: 10,48 έναντι 9,23. Η βαθμολογία και η ποιότητα εξυπηρέτησης δεν είναι το ίδιο κριτήριο.

Οι νέες επαναλήψεις εξηγούν γιατί οι μέσοι διαφέρουν από την προηγούμενη παρτίδα: Dinos στις αιχμές 185,8 αντί 191,6 χιλ. Διατηρήθηκαν και οι χαμηλές επιδόσεις, όπως η 135,5 χιλ. του Dinos στις αιχμές και η 159,4 χιλ. του Dinos+Transfer V.2 στην αλλαγή συνθηκών.

Αμφότερα παραμένουν κάτω από τις προηγουμένως μετρημένες αναφορές του ρυθμισμένου ελεγκτή: 201,9 χιλ. στις αιχμές και 225,0 χιλ. στην αλλαγή συνθηκών. Οι αναφορές προέρχονται από τα ίδια δεδομένα ανάπτυξης, όχι από νέο ανεξάρτητο τελικό έλεγχο.
\src{.lab/\allowbreak j450/\allowbreak T-\allowbreak summary.json; .lab/\allowbreak j438/\allowbreak summary.json; code-\allowbreak poisson-\allowbreak spiky/\allowbreak tuning/\allowbreak results/\allowbreak spiky-\allowbreak stage1/\allowbreak spiky1-\allowbreak P0-\allowbreak reference-\allowbreak off-\allowbreak s16000.json}
\subsection*{Τι ελέγχθηκε και τι δεν αποδεικνύεται}
Η αβεβαιότητα υπολογίστηκε με 10.000 επαναδειγματοληψίες, κρατώντας κοινές τις επιλογές διαδρομών αιτημάτων και αντιστοιχίζοντας τις επαναλήψεις εκπαίδευσης ανά κοινό αριθμό τυχαίας αρχικοποίησης. Αυτό δεν σημαίνει ταυτόσημη πορεία μάθησης. Η λεπτομερής αριθμητική σύνοψη περιλαμβάνει επίσης κεντρικό μέσο χωρίς τα άκρα, πιθανότητα καλύτερης επίδοσης, πόρους, καθυστερήσεις και απορρίψεις. Οι πλήρεις ορισμοί διατηρούνται στα συνοδευτικά αρχεία. Ο συνολικός χρόνος επεξεργαστή ήταν μικρότερος για το Dinos+Transfer V.2: περίπου 29.329 έναντι 34.480 δευτερολέπτων στις αιχμές και 34.173 έναντι 37.152 στην αλλαγή συνθηκών. Αυτό δεν αποδεικνύει ταχύτερη μάθηση.

Αυτά είναι \textbf{διερευνητικά αποτελέσματα}. Οι συνταγές είχαν ήδη επιλεγεί με δεδομένα ανάπτυξης. Δεν τεκμηριώνουν ακόμη γενίκευση σε άγνωστα τελικά δεδομένα ούτε καθαρό αιτιώδες όφελος από την προεκπαίδευση, επειδή διαφέρουν και άλλες ρυθμίσεις των δύο μεθόδων. Η εκπαίδευση απευθείας στο αλλαγμένο φορτίο εδώ είναι διαφορετικό ερώτημα από τη δοκιμή χωρίς νέα εκπαίδευση της προηγούμενης σελίδας.
\src{.lab/\allowbreak j445/\allowbreak FREEZE.md: Ανάλυση προεγγεγραμμένη; .lab/\allowbreak j450/\allowbreak T-\allowbreak summary.json; .lab/\allowbreak j450/\allowbreak summarize\_t.py}
\page{8. Τι έχει επιτευχθεί και τι απομένει πριν από τις 24/11}
\begin{tabular}{p{4.0cm}p{5.5cm}p{5.5cm}}\toprule
Στόχος & Τεκμηριωμένο σήμερα & Κενό / επόμενο βήμα\\\midrule
Η επιθυμητή σειρά μοντέλων & Υπάρχουν πλήρεις συγκρίσεις και τελικός έλεγχος στο βασικό φορτίο. & Δεν επιτεύχθηκε συνολικά. Το Dinos+Transfer V.2 δεν είναι σταθερά πρώτο· ο ρυθμισμένος ελεγκτής παραμένει ισχυρότερος σε μέση βαθμολογία.\\
Κέρδος έναντι απλών μεθόδων & Οι παλιοί τελικοί έλεγχοι στηρίζουν υπεροχή έναντι DERP. Στο νέο φορτίο στηρίζεται υπεροχή έναντι του προσαρμοσμένου DRe-SCale. & Να οριστεί με ακρίβεια ο περιορισμένος ισχυρισμός που αντέχει στις μετρήσεις. Όχι γενική υπεροχή όλων των παραλλαγών.\\
Σταθερότητα και προσαρμογή & Υπάρχουν δέκα επαναλήψεις για τις πρόσφατες συγκρίσεις και δοκιμές αλλαγής συνθηκών. & Νέος ανεξάρτητος τελικός έλεγχος για αιχμές/αλλαγή συνθηκών δεν έχει γίνει. Απαιτεί ξεχωριστό εγκεκριμένο σχέδιο.\\
Ταχύτερη μάθηση, μικρότερη χρήση πόρων & Υπάρχουν καμπύλες ανάπτυξης και μετρημένα μεγέθη κλασικών δικτύων. & Συνολική αποτίμηση προεκπαίδευσης, χρόνου, μνήμης και δίσκου με ισότιμα μικρά κλασικά μοντέλα δεν έχει ολοκληρωθεί.\\
Υβριδική κβαντική συνεισφορά & Έχουν καταγραφεί προτάσεις σχεδιασμού. & Σε παύση· δεν υπάρχουν τελικά κβαντικά αποτελέσματα. Χρειάζεται δική σου απόφαση για την επανέναρξη και το εύρος.\\
Πραγματικό Kubernetes/Redis & Έγιναν αρχικές απόπειρες βαθμονόμησης. & Σε παύση· δεν υπάρχει ολοκληρωμένη αξιολόγηση των τελικών μοντέλων.\\
Άρθρο και περίληψη συνεδρίου & Υπάρχουν σημειώσεις, σχετική βιβλιογραφία και αναφορά γραφικών. & Ολοκληρωμένο, εγκεκριμένο κείμενο και υποβολή: δεν τεκμηριώνονται.\\\bottomrule
\end{tabular}
\src{docs/\allowbreak STATE.md §§1-\allowbreak 5; docs/\allowbreak PLAN.md §1; docs/\allowbreak paper\_notes.md §§3-\allowbreak 7; docs/\allowbreak quantum\_design\_v1.md; docs/\allowbreak EXPERIMENTS.md; M-\allowbreak 20261109-\allowbreak 040; .lab/\allowbreak j450/\allowbreak T-\allowbreak summary.json}
\subsection*{Προτεινόμενη σειρά δουλειάς μέχρι την προθεσμία}
\begin{enumerate}
\item Να συμφωνηθεί το μήνυμα της δημοσίευσης με βάση τα επιβεβαιωμένα ευρήματα και τις αποτυχίες της επιθυμητής κατάταξης. Η αναφορά δεν προσθέτει νέα έγκριση πειραμάτων.
\item Να ελεγχθεί ανεξάρτητα η τελευταία ανάλυση και να αποφασιστεί αν αξίζει νέος τελικός έλεγχος στις αιχμές. Τα παλιά τελικά δεδομένα παραμένουν κλειστά για νέα επιλογή ρυθμίσεων.
\item Να αποφασιστεί αν ο στόχος παραμένει κβαντικό άρθρο σε αυτή την προθεσμία. Η εκκρεμής κβαντική συνεισφορά είναι ουσιαστικό κενό, όχι εργασία που θεωρείται ήδη ολοκληρωμένη.
\item Να γραφτούν περίληψη, μέθοδος και αποτελέσματα παράλληλα: κοινές συνθήκες σύγκρισης, όλες οι αποτυχίες, κόστος προεκπαίδευσης, όρια προσομοιωτή και σαφής διαχωρισμός ανάπτυξης/τελικού ελέγχου.
\end{enumerate}
\src{docs/\allowbreak STATE.md §1 και §4; docs/\allowbreak researcher\_requirements.md A1/\allowbreak A7, B7-\allowbreak B8; docs/\allowbreak paper\_notes.md §6; .lab/\allowbreak j445/\allowbreak FREEZE.md}
\textbf{Αβεβαιότητα χρονοδιαγράμματος:} η 24/11/2026 είναι η καταγεγραμμένη εσωτερική προθεσμία περίληψης για CCGrid. Δεν επαληθεύτηκε εξωτερικά εδώ ούτε υπάρχει τεκμηριωμένη εγγύηση ότι το σύνολο κβαντικών και πραγματικών πειραμάτων προλαβαίνει. Η διάρκεια που απομένει για αυτά είναι \textbf{άγνωστη}.
\src{docs/\allowbreak STATE.md §1; docs/\allowbreak PLAN.md §1}
\end{document}
